% GENERATED FILE. Edit canonical lessons, then run npm run generate:books.
% canonical-source-hash: fnv1a64:154589f72611b791
% canonical-lessons: BN-C425-megh, BN-C425-borsha, BN-C425-musholdhare, BN-C425-bojropat, BN-C425-chhaoya

\chapter{Clouds and Monsoon Rain}
\label{ch:bn-clouds-and-monsoon-rain}
\hlchaptermodality{425}

\begin{chapteropening}
\textbf{By the end of this chapter:} \emph{I can talk about clouds, the monsoon, a downpour and a lightning strike.}

Complete the last lesson of chapter 425: I can talk about clouds, the monsoon, a downpour and a lightning strike.
\end{chapteropening}

\section[\texorpdfstring{\bn{মেঘ} (megh)}{মেঘ (megh)}]{\bn{মেঘ} (megh) — a cloud}

\label{lesson:BN-C425-megh}



\begin{quote}
Before the new one: say the Bengali for demand (for goods), then the Bengali for to fluctuate.
\end{quote}

\subsection*{You'll want to know: \bn{মেঘ}}
\textbf{\bn{মেঘ}} — \emph{megh} — \textquotedblleft{}a cloud\textquotedblright{}.

\begin{quote}
\bn{আকাশে} \bn{কালো} \bn{মেঘ} \bn{দেখে} \bn{আমরা} \bn{ছাতা} \bn{নিলাম।}
\end{quote}

\begin{quote}
\emph{ākāshe kālo megh dekhe āmrā chhātā nilām.}
\end{quote}

\begin{quote}
Seeing dark clouds in the sky, we took umbrellas.
\end{quote}

Thunder is said to be the clouds calling: \bn{মেঘ} \bn{ডাকে}.

\begin{cousinweb}
\textbf{\bn{মেঘ}} (\emph{megh}) is Sanskrit \emph{megha}, \textquotedblleft{}a cloud\textquotedblright{} (Monier-Williams), a word Monier-Williams links with \emph{mih}, to sprinkle: the cloud is the sprinkler.
\end{cousinweb}

\begin{grammarlens}[title={What you've built}]
One more word for getting around.
\end{grammarlens}

\subsection*{Your turn}
\begin{itemize}
\raggedright
  \item \emph{Say it:} \emph{megh}
  \item \emph{Say it:} \emph{megh}, once more
  \item \emph{Say it:} \emph{ākāshe kālo megh dekhe āmrā chhātā nilām.} — Seeing dark clouds in the sky, we took umbrellas.
  \item \emph{Recall:} say the Bengali for demand (for goods), then the Bengali for to fluctuate, then say \emph{megh} again
\end{itemize}

\begin{tcolorbox}[breakable,colback=teal!4,colframe=teal!35!black,title={Before you move on}]
What does \bn{মেঘ} mean? (\textquotedblleft{}A cloud\textquotedblright{}.) Say it once more.
\end{tcolorbox}
\section[\texorpdfstring{\bn{বর্ষা} (bôrshā)}{বর্ষা (bôrshā)}]{\bn{বর্ষা} (bôrshā) — the monsoon; the rainy season}

\label{lesson:BN-C425-borsha}



\begin{quote}
Before the new one: say the Bengali for to fluctuate, then the Bengali for a cloud.
\end{quote}

\subsection*{You'll want to know: \bn{বর্ষা}}
\textbf{\bn{বর্ষা}} — \emph{bôrshā} — \textquotedblleft{}the monsoon; the rainy season\textquotedblright{}.

\begin{quote}
\bn{বর্ষায়} \bn{গ্রামের} \bn{রাস্তা} \bn{কাদায়} \bn{ভরে} \bn{যায়।}
\end{quote}

\begin{quote}
\emph{bôrshāy grāmer rāstā kādāy bhôre jāy.}
\end{quote}

\begin{quote}
In the monsoon the village roads fill up with mud.
\end{quote}

\bn{বর্ষায়}, in the rains; \bn{বর্ষাকাল}, the rainy season.

\begin{cousinweb}
\textbf{\bn{বর্ষা}} (\emph{bôrshā}) is Sanskrit \emph{varṣā}, \textquotedblleft{}the rains, the rainy season\textquotedblright{} (Monier-Williams), from \emph{vṛṣ}, to rain, the same root as \bn{বৃষ্টি}, rain.
\end{cousinweb}

\begin{grammarlens}[title={What you've built}]
One more word for getting around.
\end{grammarlens}

\subsection*{Your turn}
\begin{itemize}
\raggedright
  \item \emph{Say it:} \emph{bôrshā}
  \item \emph{Say it:} \emph{bôrshā}, once more
  \item \emph{Say it:} \emph{bôrshāy grāmer rāstā kādāy bhôre jāy.} — In the monsoon the village roads fill up with mud.
  \item \emph{Recall:} say the Bengali for to fluctuate, then the Bengali for a cloud, then say \emph{bôrshā} again
\end{itemize}

\begin{tcolorbox}[breakable,colback=teal!4,colframe=teal!35!black,title={Before you move on}]
What does \bn{বর্ষা} mean? (\textquotedblleft{}The monsoon; the rainy season\textquotedblright{}.) Say it once more.
\end{tcolorbox}
\section[\texorpdfstring{\bn{মুষলধারে} (mushôldhāre)}{মুষলধারে (mushôldhāre)}]{\bn{মুষলধারে} (mushôldhāre) — in torrents (of rain)}

\label{lesson:BN-C425-musholdhare}



\begin{quote}
Before the new one: say the Bengali for a cloud, then the Bengali for the monsoon.
\end{quote}

\subsection*{You'll want to know: \bn{মুষলধারে}}
\textbf{\bn{মুষলধারে}} — \emph{mushôldhāre} — \textquotedblleft{}in torrents (of rain)\textquotedblright{}.

\begin{quote}
\bn{সারারাত} \bn{মুষলধারে} \bn{বৃষ্টি} \bn{হল।}
\end{quote}

\begin{quote}
\emph{sārārāt mushôldhāre brishṭi hôlô.}
\end{quote}

\begin{quote}
It poured all night.
\end{quote}

Always with rain: \bn{মুষলধারে} \bn{বৃষ্টি}, rain in torrents.

\begin{cousinweb}
\textbf{\bn{মুষলধারে}} (\emph{mushôldhāre}) joins Sanskrit \emph{musala}, \textquotedblleft{}a pestle\textquotedblright{}, and \emph{dhārā}, \textquotedblleft{}a stream, a jet\textquotedblright{} (Monier-Williams for both), with the locative ending: rain falling in streams as thick as pestles.
\end{cousinweb}

\begin{grammarlens}[title={What you've built}]
One more word for getting around.
\end{grammarlens}

\subsection*{Your turn}
\begin{itemize}
\raggedright
  \item \emph{Say it:} \emph{mushôldhāre}
  \item \emph{Say it:} \emph{mushôldhāre}, once more
  \item \emph{Say it:} \emph{sārārāt mushôldhāre brishṭi hôlô.} — It poured all night.
  \item \emph{Recall:} say the Bengali for a cloud, then the Bengali for the monsoon, then say \emph{mushôldhāre} again
\end{itemize}

\begin{tcolorbox}[breakable,colback=teal!4,colframe=teal!35!black,title={Before you move on}]
What does \bn{মুষলধারে} mean? (\textquotedblleft{}In torrents (of rain)\textquotedblright{}.) Say it once more.
\end{tcolorbox}
\section[\texorpdfstring{\bn{বজ্রপাত} (bôjrôpāt)}{বজ্রপাত (bôjrôpāt)}]{\bn{বজ্রপাত} (bôjrôpāt) — a lightning strike; a thunderbolt}

\label{lesson:BN-C425-bojropat}



\begin{quote}
Before the new one: say the Bengali for the monsoon, then the Bengali for in torrents.
\end{quote}

\subsection*{You'll want to know: \bn{বজ্রপাত}}
\textbf{\bn{বজ্রপাত}} — \emph{bôjrôpāt} — \textquotedblleft{}a lightning strike; a thunderbolt\textquotedblright{}.

\begin{quote}
\bn{ঝড়ের} \bn{রাতে} \bn{গ্রামের} \bn{একটা} \bn{উঁচু} \bn{গাছে} \bn{বজ্রপাত} \bn{হল।}
\end{quote}

\begin{quote}
\emph{jhôṛer rāte grāmer ekṭā ũchu gāchhe bôjrôpāt hôlô.}
\end{quote}

\begin{quote}
On the stormy night lightning struck a tall tree in the village.
\end{quote}

\bn{বজ্র} is thunder itself; \bn{বজ্রপাত} is the bolt that falls.

\begin{cousinweb}
\textbf{\bn{বজ্রপাত}} (\emph{bôjrôpāt}) is Sanskrit \emph{vajra-pāta}, \textquotedblleft{}the fall of a thunderbolt, a stroke of lightning\textquotedblright{} (Monier-Williams), from \textbf{\bn{বজ্র}}, thunder, and \emph{pāta}, a fall.
\end{cousinweb}

\begin{grammarlens}[title={What you've built}]
One more word for getting around.
\end{grammarlens}

\subsection*{Your turn}
\begin{itemize}
\raggedright
  \item \emph{Say it:} \emph{bôjrôpāt}
  \item \emph{Say it:} \emph{bôjrôpāt}, once more
  \item \emph{Say it:} \emph{jhôṛer rāte grāmer ekṭā ũchu gāchhe bôjrôpāt hôlô.} — On the stormy night lightning struck a tall tree in the village.
  \item \emph{Recall:} say the Bengali for the monsoon, then the Bengali for in torrents, then say \emph{bôjrôpāt} again
\end{itemize}

\begin{tcolorbox}[breakable,colback=teal!4,colframe=teal!35!black,title={Before you move on}]
What does \bn{বজ্রপাত} mean? (\textquotedblleft{}A lightning strike; a thunderbolt\textquotedblright{}.) Say it once more.
\end{tcolorbox}
\section[\texorpdfstring{\bn{ছাওয়া} (chhāoyā)}{ছাওয়া (chhāoyā)}]{\bn{ছাওয়া} (chhāoyā) — to spread over; to cover over (of clouds); to thatch}

\label{lesson:BN-C425-chhaoya}



\begin{quote}
Before the new one: say the Bengali for in torrents, then the Bengali for a lightning strike.
\end{quote}

\subsection*{You'll want to know: \bn{ছাওয়া}}
\textbf{\bn{ছাওয়া}} — \emph{chhāoyā} — \textquotedblleft{}to spread over; to cover over (of clouds); to thatch\textquotedblright{}.

\begin{quote}
\bn{দুপুরে} \bn{আকাশ} \bn{মেঘে} \bn{ছেয়ে} \bn{গেল।}
\end{quote}

\begin{quote}
\emph{dupure ākāsh meghe chheye gelô.}
\end{quote}

\begin{quote}
At midday clouds spread across the sky.
\end{quote}

Past (he/she …): \bn{সেই} \bn{আকাশ} \bn{মেঘে} \bn{ছেয়ে} \bn{গেল}, \emph{sei ākāsh meghe chheye gelô}.

Mostly in the -\bn{এ} form with \bn{যাওয়া}: \bn{ছেয়ে} \bn{যাওয়া}, to be covered all over. Also to thatch a roof.

\begin{cousinweb}
\textbf{\bn{ছাওয়া}} (\emph{chhāoyā}) is inherited: CDIAL lists Bengali \emph{chāuyā}, \textquotedblleft{}to cover, to thatch\textquotedblright{}, under Sanskrit \emph{chādáyati}, \textquotedblleft{}covers, veils\textquotedblright{} (CDIAL 5018), the root of \bn{ছই}, a boat's hood, as well.
\end{cousinweb}

\begin{grammarlens}[title={What you've built}]
One more word for getting around.
\end{grammarlens}

\subsection*{Your turn}
\begin{itemize}
\raggedright
  \item \emph{Say it:} \emph{chhāoyā}
  \item \emph{Say it:} \emph{chhāoyā}, once more
  \item \emph{Say it:} \emph{dupure ākāsh meghe chheye gelô.} — At midday clouds spread across the sky.
  \item \emph{Recall:} say the Bengali for in torrents, then the Bengali for a lightning strike, then say \emph{chhāoyā} again
\end{itemize}

\begin{tcolorbox}[breakable,colback=teal!4,colframe=teal!35!black,title={Before you move on}]
What does \bn{ছাওয়া} mean? (\textquotedblleft{}To spread over; to cover over (of clouds); to thatch\textquotedblright{}.) Say it once more.
\end{tcolorbox}
